% =====================================================================
% Discussion section draft — Mate-choice copying update meta-analysis
% Drafted to slot into the getwriting manuscript (biblatex citations).
%
% NOTE ON CITATIONS: keys below were taken from references.bib in the
% supplement repo. If the manuscript .bib uses different keys, remap:
%   jonesMechanismsSocialInfluence2019      = Jones & DuVal 2019  [5]
%   daviesMetaanalysisFactorsInfluencing2020 = Davies et al. 2020  [6]
%   vakirtzisMateChoiceCopying2011          = Vakirtzis 2011       [2]
%   nakagawaMethodsTestingPublication2022   = Nakagawa et al. 2022 [31]
%   yangRobustPointVariance2024             = Yang et al. 2024     [34]
%   belkinaMateChoiceCopying2021            = Belkina et al. 2021  [36]
%   cirinoRobustMatePreferences2023         = Cirino et al. 2023   [38]
%   pusiakRelativeSexualAttractiveness2020  = Pusiak et al. 2020   [60]
%   nobelNoEvidenceMate2023                 = Nobel et al. 2023    [58]
%   <pollo2025key>                          = Pollo et al. 2025    [9]
% =====================================================================

\section{Discussion}

By updating the datasets of the two source meta-analyses and re-analysing them
with a unified, phylogenetically informed multilevel framework, we found that
the evidence for mate-choice copying in non-human animals remains positive and
statistically robust, but is more modest than the original syntheses implied.
On the combined data, the pooled odds ratio was 1.81 (95\% CI: 1.35 to 2.42)
and the pooled Hedges' \textit{g} was 0.445 (95\% CI: 0.241 to 0.649): observers
exposed to social information were reliably more likely to favour the
demonstrated mate than observers that were not. Both estimates, however, sit
appreciably below those reported by \textcite{jonesMechanismsSocialInfluence2019}
(OR = 2.71) and \textcite{daviesMetaanalysisFactorsInfluencing2020}
(Hedges' \textit{d} = 0.58). The attenuation is most evident when the new
evidence is considered on its own: the new-extraction means
(Hedges' \textit{g} = 0.206, 95\% CI: $-$0.017 to 0.431; OR = 1.46, 95\% CI:
0.968 to 2.201) are roughly half the magnitude of the originals and have
confidence intervals that span the null. The qualitative conclusion of the two
foundational meta-analyses therefore survives an expanded, multilingual, and
grey-literature--inclusive search, but the magnitude of mate-choice copying
appears to have been somewhat overestimated---a pattern that is itself
informative and that we explore below.

\subsection{Knowledge gaps and future opportunities}

The most conspicuous feature of the updated evidence base is its persistent and
largely unexplained heterogeneity. Total heterogeneity exceeded 74\% in every
model, yet very little of it was attributable to differences between studies
(study-level $I^2 < 8\%$) or to shared ancestry (phylogenetic $I^2 < 7\%$).
Instead, variation accumulated at the within-study (observation) and
non-phylogenetic species levels. In other words, the disagreement among
mate-choice copying estimates is not primarily a matter of some laboratories
producing systematically larger effects than others, nor of copying being
concentrated in particular branches of the tree of life; it arises within
studies and species, where the experimental and ecological detail lives. This
points to genuine biological and methodological moderators that our intercept
models could not capture, and identifying them is, in our view, the single most
valuable direction for future work. We collected several candidate moderators
during extraction---the copying mechanism (generalised versus individual), the
modality of the social cue (visual, chemical, or acoustic), the virginity of the
observer, and the relative age of the demonstrator---and we encourage future
syntheses, ideally pre-registered, to test these explicitly as the dataset
continues to grow.

A second gap concerns taxonomic breadth. Although our combined sample spanned 33
species across arthropods, fish, and other vertebrates, the evidence remains
heavily concentrated in two model genera: \textit{Drosophila} and
\textit{Poecilia} together contributed 66\% of all effect sizes. Reassuringly,
the uni-moderator analyses returned positive means in all three taxonomic groups
on both scales, and we found only weak statistical support for among-group
differences (Hedges' \textit{g}: $Q_M = 3.06$, $p = 0.06$; log odds ratio:
$Q_M = 1.17$, $p = 0.32$). Nonetheless, the point estimates hint at a gradient
worth pursuing---copying appeared strongest in birds and mammals
(Hedges' \textit{g} = 0.71) and weakest, though still positive, in arthropods
(Hedges' \textit{g} = 0.24)---and the confidence we can place in any taxonomic
contrast is limited by how few clades are well sampled. Amphibians, reptiles,
non-poeciliid fishes, and a broader range of birds and mammals are effectively
absent. Expanding the phylogenetic spread of the evidence would not only test
the generality of mate-choice copying but would also sharpen the power of
phylogenetic comparative tests that, at present, have little signal to work
with.

The third and most consequential gap is the strength of the small-study and
selective-reporting effects we detected. The Egger-type multilevel slopes on
sampling variance were positive and clearly significant on both scales
(Hedges' \textit{g}: $\beta_{\sqrt{v_i}} = 4.36$, 95\% CI: 3.22 to 5.49; odds
ratio: $\beta_{\sqrt{v_i}} = 2.11$, 95\% CI: 1.23 to 2.99), indicating that
smaller, less precise studies reported disproportionately large effects. The two
correction frameworks bracket the problem instructively. Under the bias-corrected
intercept of \textcite{nakagawaMethodsTestingPublication2022}, the adjusted
combined means collapsed towards---and for Hedges' \textit{g} marginally below---
zero (Hedges' \textit{g} = $-$0.21; log odds ratio = 0.10), whereas the
bias-robust estimator of \textcite{yangRobustPointVariance2024} retained a
positive but reduced signal (Hedges' \textit{g} = 0.34; log odds ratio = 0.57).
The honest interpretation is that the true effect very likely lies below the
naive pooled estimate, and that we cannot, on present evidence, confidently
exclude a much smaller effect. This is not a reason to dismiss mate-choice
copying, but it is a strong argument for the field to adopt registered reports
and to routinely publish null and small-sample results, so that future updates
are anchored on a less censored literature.

Related to this is a partial decline effect. For the binary outcomes we found
that effect-size magnitude decreased with publication year
($\beta_{\mathrm{year}} = -0.028$, 95\% CI: $-0.048$ to $-0.007$), whereas for the
continuous outcomes the trend was negligible ($\beta_{\mathrm{year}} = -0.002$,
95\% CI: $-0.014$ to $0.01$). A decline that appears on one scale but not the
other should be interpreted cautiously, but it is consistent with a literature
in which striking early demonstrations are gradually tempered by more rigorous
and better-powered designs. Several of the most recent studies in our sample
report weak or null copying---for example in \textit{Drosophila}
\parencite{belkinaMateChoiceCopying2021,cirinoRobustMatePreferences2023}, in
guppies \parencite{pusiakRelativeSexualAttractiveness2020}, and the absence of
copying in zebrafish \parencite{nobelNoEvidenceMate2023}---and these temper the
early enthusiasm that the phenomenon is both ubiquitous and strong.

Finally, our results expose a methodological fragmentation that constrains
synthesis itself. Primary studies quantify copying in incompatible ways---from
preference-zone time to social-learning indices to discrete copulation
choices---and the two source meta-analyses reported their results in different
metrics, forcing us to convert between standardised mean differences and log
odds ratios under an approximate logistic-to-normal assumption. The
cross-dataset validation underscored the cost of this fragmentation: seven
studies disagreed between the original meta-analyses by more than 0.5
Hedges' \textit{g} units. Encouragingly, our headline estimates were robust to
how these contentious values were handled---leave-one-out, leave-one-species-out,
and the conservative ``leave-out'' and ``corrected'' re-analyses all left the
sign and significance of both syntheses unchanged, with the largest shift in any
pooled mean being only 0.04--0.07 units. Even the removal of \textit{Poecilia
reticulata}, the most-studied and strongest-copying species, merely nudged the
mean towards zero without overturning it. The broader lesson, echoing
\textcite{daviesMetaanalysisFactorsInfluencing2020} and recent calls for
verification in our field, is that the community would benefit from minimum
reporting standards---group means, standard deviations, sample sizes, and raw
choice counts---so that future updates rest on directly comparable, fully
re-extractable data.

\subsection{Conclusion}

Mate-choice copying is a real and taxonomically widespread behaviour: after
incorporating seven additional years of evidence and analysing it with
contemporary phylogenetic multilevel methods, we still recover a positive and
significant overall effect on both the odds-ratio and standardised-mean-difference
scales. At the same time, the updated picture is more sober than the original
meta-analyses suggested. The pooled effect is moderate at best, it is surrounded
by substantial heterogeneity that our models could not attribute to study,
taxon, or phylogeny, and it bears clear fingerprints of small-study and
selective-reporting bias that, under the more conservative correction, draw the
estimate close to zero. The qualitative claims of
\textcite{jonesMechanismsSocialInfluence2019} and
\textcite{daviesMetaanalysisFactorsInfluencing2020} hold; their magnitudes were
probably optimistic. Realising the promise of mate-choice copying as a driver of
sexual selection and cultural evolution will require broader taxonomic sampling,
mechanism-focused and pre-registered tests of the moderators that drive its
considerable heterogeneity, and a publication culture that records nulls as
faithfully as it records discoveries. Beyond its specific findings, we offer this
update as a template for how meta-analyses in ecology and evolution can be
revisited---combining new evidence with the original data, harmonising
incompatible effect-size metrics, and stress-testing conclusions against
publication bias and contentious extractions.
